\documentclass[10pt,a4paper]{article}

% ----------------- Paquetes -------------------%
\usepackage[T1]{fontenc}
\usepackage[utf8]{inputenc}
\usepackage[a4paper,margin=2.5cm]{geometry}
\usepackage{mathptmx}             
\usepackage{changepage} 
\usepackage{setspace}
\usepackage[spanish]{babel}
\usepackage[labelfont=bf,labelsep=space]{caption}
\usepackage{amssymb}
\usepackage{amsmath}
\usepackage{ebgaramond}
\usepackage{placeins}
\usepackage{etoolbox}
\usepackage{setspace}
\usepackage{parskip} 
\usepackage{tcolorbox}
\usepackage{needspace}
\usepackage{xparse}
\usepackage{ocgx2}
\usepackage{hyperref}
\usepackage{hypcap}


%%%%%%%%%%%%%%%%%%%%%%%%%%%%%%%%%%%%%%%%%%%%%%%%%%%%%%%%%%%%%%%%
% --- Fija primero tus valores globales "buenos" ---
\setstretch{1.2}                 % Interlineado global
\setlength{\parskip}{0.7\baselineskip}
\setlength{\parindent}{0pt}
\newlength{\GlobalParskip}
\setlength{\GlobalParskip}{\parskip}

\newcounter{eqnum}[subsection]
\renewcommand{\theeqnum}{\arabic{eqnum}}
\newcounter{alphaitem}
\newcounter{numitem}

%% ------------------- Recuadros----------------------%%
\newcommand{\puntosclave}[1]{%
  \begin{tcolorbox}[
    colframe=black,
    title={Puntos clave},
    before upper={\setlength{\parindent}{0.5cm}\setlength{\parskip}{0.5\baselineskip}}
  ]
  #1
  \end{tcolorbox}
}

\newcommand{\modificaciones}[1]{
  \begin{tcolorbox}[
    colback=white,
    colframe=red,
    title={Modificaciones a realizar},
    before upper={\setlength{\parindent}{0.5cm}\setlength{\parskip}{0.5\baselineskip}}
  ]
  #1
  \end{tcolorbox}
}
%% ------------------- Preguntas TEST----------------------%%
% Opciones: radio (single-choice)
\NewDocumentCommand{\OptRadio}{m m m}{%
  \noindent\CheckBox[radio,name=#1,value=#2,width=1.2ex,height=1.2ex]{}%
  \;\textbf{#2.}~#3\par
}

% Opciones: checkbox (multi-choice)
\NewDocumentCommand{\OptCheck}{m m}{%
  \noindent\CheckBox[name=#1,width=1.2ex,height=1.2ex,bordercolor={0 0 0}]{}%
  \;#2\par
}

% Pregunta single-choice (radio)
% Args: 1=id, 2=enunciado, 3=opciones, 4=solución+justificación, 5=imagen (o {}), 6=origen
\NewDocumentCommand{\PreguntaTestSingle}{m m m m m m}{%
  \par\medskip
  \noindent\textbf{#1.}~#2\par\smallskip

    % Imagen (si no es {})
    \IfBlankTF{#5}{}{%
      \begin{center}
        \includegraphics[width=0.75\linewidth]{#5}
      \end{center}
    }

    \begin{Form}
      #3
    \end{Form}

    \par\smallskip
    \switchocg{sol-#1}{\fbox{Mostrar / ocultar solución}}%
    \hfill{\footnotesize #6}\par\smallskip

    \begin{ocg}{Solución}{sol-#1}{0}
      \noindent #4
    \end{ocg}
  \par\medskip
}

% Pregunta multi-choice (checkboxes)
\NewDocumentCommand{\PreguntaTestMulti}{m m m m m m}{%
  \par\medskip
  \noindent\textbf{#1.}~#2\par\smallskip

    % Imagen (si no es {})
    \IfBlankTF{#5}{}{%
      \begin{center}
        \includegraphics[width=0.75\linewidth]{#5}
      \end{center}
    }
    \begin{Form}
      #3
    \end{Form}

    \par\smallskip
    \switchocg{sol-#1}{\fbox{Mostrar / ocultar solución}}%
    \hfill{\footnotesize #6}\par\smallskip

    \begin{ocg}{Solución}{sol-#1}{0}
      \noindent #4
    \end{ocg}
  \par\medskip
}

%% ------------------- Listas----------------------%%
    % --- Lista alfabética ---
    \newenvironment{la}
      {\begin{list}{\alph{alphaitem})}{%
          \usecounter{alphaitem}%
          \setlength{\leftmargin}{1cm}%
          \setlength{\parskip}{3pt}% 
          \setlength{\itemsep}{0pt}% ← espacio entre ítems
          \setlength{\parsep}{0pt}%  ← párrafos dentro de un ítem
      }}
      {\end{list}}
      
    % --- Lista numérica ---
    \newenvironment{lnum}
      {\begin{list}{\arabic{numitem}.}{%
          \usecounter{numitem}%
          \setlength{\leftmargin}{1cm}%
          \setlength{\parskip}{3pt}% 
          \setlength{\itemsep}{0pt}% ← espacio entre ítems
          \setlength{\parsep}{0pt}%  ← párrafos dentro de un ítem
      }}
      {\end{list}}
      
% ------------------- Imágenes----------------------%
    \graphicspath{{Img/}}% Rutas por defecto 
    % --- Imagen adaptativa ---
    % Registros de dimensión definidos una sola vez
    \newdimen\imgcap
    \newdimen\imgmin
    \newdimen\imgmax
    \newdimen\imgremain
    \newdimen\imgh
    \newdimen\imgneed
    
    \newcommand{\imagenfit}[3]{%
      \addvspace{10pt}% separación antes
      \par\begingroup
        % Parámetros de composición (asignaciones, no \newdimen)
        \imgcap=1\baselineskip            % alto estimado de título+fuente
        \imgmin=0.15\textheight
        \imgmax=0.15\textheight
        \imgremain=\pagegoal
        \ifdim\pagegoal=\maxdimen \imgremain=0pt\fi
        \advance\imgremain by -\pagetotal
        \advance\imgremain by -\imgcap
        % Decisión de altura destino
        \ifdim\imgremain<\imgmin
          \newpage
          \imgh=0.20\textheight
        \else
          \imgh=\imgremain
          \ifdim\imgh>\imgmax \imgh=\imgmax \fi
        \fi
        % Reservar espacio para evitar cortes del bloque completo
        \imgneed=\imgh \advance\imgneed by \imgcap
        \Needspace{\imgneed}
        % Cuerpo
        \noindent\begin{minipage}{\textwidth}
          \centering
          \captionof{figure}{\textemdash\ #2}\par\vspace{0.3em}% título arriba
          \includegraphics[width=\linewidth,height=\imgh,keepaspectratio]{\detokenize{#1}}\par
          \vspace{0.3em}\small\textit{Fuente: }#3% fuente abajo
        \end{minipage}%
      \endgroup
      \par\addvspace{10pt}%
    }

%------------ Bloque de RECUADRO ESPECIAL -------------%
    \newenvironment{specialblock}[1]{%
      \begin{quote} % crea sangría a izquierda y derecha
      \small % texto más pequeño
      \noindent\textbf{#1}\\[-0.3em] % título en negrita
      \rule{\linewidth}{0.4pt}\\[0.6em]}{
      \end{quote}}
      
%-------------------------- Portada -----------------------%
\newcommand{\oldtitle}[1]{\def\OldTitle{#1}}
\newcommand{\changerationale}[1]{\def\ChangeWhy{#1}}

\newcommand{\changebox}{%
\begin{tcolorbox}[title={Historial de cambios del tema}]
\textbf{Título anterior:} \OldTitle\par
\textbf{Motivo del cambio:} \ChangeWhy
\end{tcolorbox}
}

%-------------------------- Author Font -----------------------%
    \newcommand{\authorfont}[1]{{\fontfamily{EBGaramond-TLF}\selectfont #1}}

%-------------------------- Anexos -----------------------%
    \makeatletter
    \newcounter{anexo}
    \renewcommand{\theanexo}{\Roman{anexo}}
    \newcommand{\anexo}[1]{%
      \clearpage
      \phantomsection
      \refstepcounter{anexo}%
      \noindent\textbf{Anexo \theanexo.- }#1\par\medskip
      \addcontentsline{stc}{section}{Anexo \theanexo.- #1}%
    }
    \makeatother

%-------------------------- Lista de figuras (personal) -----------------------%
\makeatletter
\newcommand{\MyListOfFigures}{%
  \clearpage
  \phantomsection
  {\centering\bfseries\Large Índice de figuras\par}%
  \medskip
  % Entrada en nuestro índice de contenido (stc)
  \addcontentsline{stc}{section}{Índice de figuras}%
  % Recuperación del listado (sin cabecera propia)
  \begingroup
    \parindent=0pt\parskip=2pt
    \@starttoc{lof}%
  \endgroup
}
\makeatother

%--------- Establecer cuatro niveles de índice --------%
    \setcounter{tocdepth}{4}   
    \setcounter{secnumdepth}{4}% numera hasta \paragraph

%-------------------------- Bibliografía -----------------------%
\makeatletter
% Contador para las referencias
\newcounter{myref}
% Cabecera de la bibliografía, entra en el índice 'stc'
\newcommand{\MyBibliography}{%
  \clearpage
  \phantomsection
  {\centering\bfseries\Large Bibliografía y referencias\par}%
  \medskip
  \addcontentsline{stc}{section}{Bibliografía y referencias}%
  \setcounter{myref}{0}%
}
\newcommand{\MyBibItem}[4]{%
  \refstepcounter{myref}%
  \noindent[\themyref]\ #1.\ \emph{#2}. #3. #4\par
}
\makeatother


%%%%%%%%%%%%%%%%%%%%%%%%%%%%%%%%

\makeatletter

    %--------- Interlineado Global --------%
    \edef\GlobalBaseLineStretch{\baselinestretch}
    
    %--------- Espaciado de seccinones --------%
    \renewcommand\section{\@startsection{section}
      {1}{\z@}
      {2.5ex \@plus 1ex \@minus .2ex} % espacio antes
      {1.5ex \@plus .2ex}             % espacio después
      {\normalfont\Large\bfseries}    % formato del título
    }
    \renewcommand\subsection{\@startsection{subsection}{2}{\z@}
      {3ex \@plus 0.8ex \@minus .2ex}
      {1.5ex \@plus .2ex}
      {\normalfont\large\bfseries}
    }
    \renewcommand\subsubsection{\@startsection{subsubsection}{3}{\z@}
      {3ex \@plus 0.5ex \@minus .2ex}
      {1ex \@plus .2ex}
      {\normalfont\normalsize\bfseries}
    }
    \renewcommand\paragraph{\@startsection{paragraph}{4}{\z@}%
    {-2ex \@plus -0.5ex \@minus -.2ex}% espacio antes
    {0.8ex \@plus .2ex}% espacio después
    {\normalfont\normalsize\itshape}}% ← cursiva en lugar de negrita

    %--------- Ecuaciones --------%   
    % Variables globales para almacenar títulos y notas
    \gdef\leq@current@section{}%
    \gdef\leq@current@subsection{}%
    \gdef\leq@last@printed@section{}%
    \gdef\leq@last@printed@subsection{}%
    
    % Comando para añadir nota (invisible en el cuerpo)
    \newcommand{\eqnote}[1]{%
      \addcontentsline{leq}{leqnote}{#1}%
    }
    
    % Comando para ecuaciones
    \newcommand{\eqblock}[2]{%
      % Verificar si necesitamos imprimir el título de sección
      \ifx\leq@current@section\leq@last@printed@section\else
        \addcontentsline{leq}{leqsection}{\leq@current@section}%
        \global\let\leq@last@printed@section\leq@current@section
        \gdef\leq@last@printed@subsection{}% Reset subsección al cambiar sección
      \fi
      % Verificar si necesitamos imprimir el título de subsección
      \ifx\leq@current@subsection\leq@last@printed@subsection\else
        \ifx\leq@current@subsection\@empty\else
          \addcontentsline{leq}{leqsubsection}{\leq@current@subsection}%
          \global\let\leq@last@printed@subsection\leq@current@subsection
        \fi
      \fi
      % Ecuación
      \refstepcounter{eqnum}%
      \begin{equation*}
        #1\tag{E.\theeqnum}
      \end{equation*}%
      \addcontentsline{leq}{leqt}{[E.\theeqnum] -- #2}%
      \addcontentsline{leq}{leqm}{#1}%
    }
    
    %%% ----- Índice de ecuaciones ----------%%%
    % Tamaño para []
    \newcommand{\leqIndexFont}{\normalsize}
    \newcommand{\leqTitleFont}{\tiny}
    \newcommand{\leqNoteFont}{\tiny}
    % Cabecera de sección 
    \newcommand*\l@leqsection[2]{%
      \addvspace{25pt}% Mayor separación antes de nueva sección
      \noindent\leqIndexFont
      \makebox[\linewidth]{\rule{.38\linewidth}{0.4pt}\ \ \textbf{  \MakeUppercase{#1}}\ \ \rule{.38\linewidth}{0.4pt}}%
      \par\vspace{-10pt}
    }
    % Cabecera de subsección 
    \newcommand*\l@leqsubsection[2]{%
      \par\addvspace{15pt}%
      \noindent\leqIndexFont #1
      \par\vspace{-7pt}
    }
    % Título de ecuación 
    \newcommand*\l@leqt[2]{%
    \par\addvspace{10pt}%
      \par\noindent\leqTitleFont #1\par
    }
    % Miniatura de ecuación
    \newcommand*\l@leqm[2]{%
      \vspace{0.3\baselineskip}%
      \par\scriptsize\noindent\hspace*{1.5em}\(#1\)\par%
    }
    % Nota de ecuación 
    \newcommand*\l@leqnote[2]{%
      \vspace{0.2\baselineskip}%
      \par\noindent\leqNoteFont\hspace*{1.5em}\textbf{Nota.--} #1\par\vspace{0.5\baselineskip}%
    }
    % Generación del índice
    \newcommand{\mylistofequations}{%
      \clearpage
      \phantomsection
      % Título visual de la lista (ajústalo a tu gusto):
      {\centering\bfseries\Large Índice de ecuaciones\par}
      % Entrada en nuestro índice 'stc':
      \addcontentsline{stc}{section}{Índice de ecuaciones}%
      % Llamada a tu lista real (usa el nombre exacto de tu macro):
      \@starttoc{leq}%
    }
    
    %--------- Captura de títulos de sección --------%
    \let\leq@original@section\section
    \let\leq@original@subsection\subsection
    % Flag para saber si estamos en una sección asterisco
    \newif\ifLeqStarSection
    \LeqStarSectionfalse
    
    % Redefinir \section CON CONTROL DE ASTERISCO
    \renewcommand{\section}{%
      \@ifstar{\leq@section@star}{\@dblarg\leq@section@nostar}%
    }
    \def\leq@section@star#1{%
      \global\LeqStarSectiontrue
      \gdef\leq@current@section{#1}%
      \gdef\leq@current@subsection{}%
      \leq@original@section*{#1}%
      \global\LeqStarSectionfalse
    }
    \def\leq@section@nostar[#1]#2{%
      \global\LeqStarSectionfalse
      \gdef\leq@current@section{#2}%
      \gdef\leq@current@subsection{}%
      \leq@original@section[#1]{#2}%
    }
    % Redefinir \subsection
    \renewcommand{\subsection}{%
      \@ifstar{\leq@subsection@star}{\@dblarg\leq@subsection@nostar}%
    }
    \def\leq@subsection@star#1{%
      \global\LeqStarSectiontrue
      \gdef\leq@current@subsection{#1}%
      \leq@original@subsection*{#1}%
      \global\LeqStarSectionfalse
    }
    \def\leq@subsection@nostar[#1]#2{%
      \global\LeqStarSectionfalse
      \gdef\leq@current@subsection{#2}%
      \leq@original@subsection[#1]{#2}%
    }

    % ----- Restauración de valores Globales de estilo ----------%
    % Flags
    \newif\ifInFootnote
    \InFootnotefalse
    \pretocmd{\@footnotetext}{\InFootnotetrue}{}{}
    \apptocmd{\@footnotetext}{\InFootnotefalse}{}{}
    % Profundidad de listas (soporta anidadas)
    \newcount\MyListDepth \MyListDepth=0
    \AtBeginEnvironment{la}{\advance\MyListDepth by 1}
    \AfterEndEnvironment{la}{\advance\MyListDepth by -1}
    \AtBeginEnvironment{lnum}{\advance\MyListDepth by 1}
    \AfterEndEnvironment{lnum}{\advance\MyListDepth by -1}
    \AtBeginEnvironment{itemize}{\advance\MyListDepth by 1}
    \AfterEndEnvironment{itemize}{\advance\MyListDepth by -1}
    \AtBeginEnvironment{enumerate}{\advance\MyListDepth by 1}
    \AfterEndEnvironment{enumerate}{\advance\MyListDepth by -1}
    % Restauración diferida: hazlo al INICIO del próximo párrafo real
    \newif\if@pendingRestore
    \@pendingRestorefalse
    \newcommand{\RequestRestore}{\global\@pendingRestoretrue}
    \everypar{%
      \if@pendingRestore
        \global\@pendingRestorefalse
        \ifInFootnote\else
          \ifnum\MyListDepth=0
            \setlength{\parskip}{\GlobalParskip}%
            \setstretch{\GlobalBaseLineStretch}%
          \fi
        \fi
      \fi
    }
    % Al final de bloques:
    \AfterEndEnvironment{equation}{\RequestRestore}
    \AfterEndEnvironment{equation*}{\RequestRestore}
    \AfterEndEnvironment{la}{\RequestRestore}
    \AfterEndEnvironment{lnum}{\RequestRestore}
    % Al final de macros:
    \apptocmd{\eqblock}{\RequestRestore}{}{}
    \apptocmd{\imagenfit}{\RequestRestore}{}{}
    
\makeatother

% ===================== ToC propio con 4 niveles (único bloque) =====================
\makeatletter

% --- Escritores: añadimos una línea al .stc en cada cabecera numerada ---
% Guardamos las versiones actuales para llamar al comportamiento normal
\let\MyTOC@orig@section\section
\let\MyTOC@orig@subsection\subsection
\let\MyTOC@orig@subsubsection\subsubsection
\let\MyTOC@orig@paragraph\paragraph

% SECTION
\renewcommand{\section}{%
  \@ifstar{\MyTOC@section@star}{\@dblarg\MyTOC@section@nostar}%
}
\def\MyTOC@section@star#1{%
  \MyTOC@orig@section*{#1}% sin entrada al .stc
}
\def\MyTOC@section@nostar[#1]#2{%
  \MyTOC@orig@section[{#1}]{#2}% comportamiento normal (escribe al .toc)
  % Escribimos también nuestra línea al .stc (texto con \numberline)
  \addcontentsline{stc}{section}{\protect\numberline{\thesection}#2}%
}

% SUBSECTION
\renewcommand{\subsection}{%
  \@ifstar{\MyTOC@subsection@star}{\@dblarg\MyTOC@subsection@nostar}%
}
\def\MyTOC@subsection@star#1{%
  \MyTOC@orig@subsection*{#1}%
}
\def\MyTOC@subsection@nostar[#1]#2{%
  \MyTOC@orig@subsection[{#1}]{#2}%
  \addcontentsline{stc}{subsection}{\protect\numberline{\thesubsection}#2}%
}

% SUBSUBSECTION
\renewcommand{\subsubsection}{%
  \@ifstar{\MyTOC@subsubsection@star}{\@dblarg\MyTOC@subsubsection@nostar}%
}
\def\MyTOC@subsubsection@star#1{%
  \MyTOC@orig@subsubsection*{#1}%
}
\def\MyTOC@subsubsection@nostar[#1]#2{%
  \MyTOC@orig@subsubsection[{#1}]{#2}%
  \addcontentsline{stc}{subsubsection}{\protect\numberline{\thesubsubsection}#2}%
}

% PARAGRAPH
\renewcommand{\paragraph}{%
  \@ifstar{\MyTOC@paragraph@star}{\@dblarg\MyTOC@paragraph@nostar}%
}
\def\MyTOC@paragraph@star#1{%
  \MyTOC@orig@paragraph*{#1}%
}
\def\MyTOC@paragraph@nostar[#1]#2{%
  \MyTOC@orig@paragraph[{#1}]{#2}%
  % Si no numeras los \paragraph, cambia \theparagraph por texto plano sin \numberline
  \addcontentsline{stc}{paragraph}{\protect\numberline{\theparagraph}#2}%
}

% --- Impresor del índice propio (lee el .stc) ---
\newcommand{\MyTOC}{%
  \par\bigskip
  {\centering\bfseries\Large Índice de contenido\par}%
  \medskip
  \begingroup
    \parindent=0pt\parskip=2pt
    \@starttoc{stc}%
  \endgroup
}

\makeatother
% ===================== fin ToC propio =====================


%%%%%%%%%%%%%%


% --- Datos del tema ---
\title{Tema 3.A.26 -- Extensiones a la Teoría del Mercado de Trabajo\\
\large Desempleo friccional. La Curva de Beveridge. El modelo de búsqueda y emparejamiento de Diamond, Mortensen y Pissarides. Costes de ajuste y dinámica de la demanda de trabajo.}
\author{Marcelo Ortega}
\date{Noviembre 2025}
\newcommand{\TemaCode}{3A26}

\oldtitle{Tema 3.A.19. Extensiones de las teorías de oferta y demanda de trabajo: información, búsqueda, costes de ajuste y dinámica de la demanda de trabajo}
\changerationale{Se clarifica que la problemática central del tema es el desempleo friccional.}


\usepackage{soul}\sethlcolor{yellow}% resaltado amarillo: \hl{texto} (Panel TCEE, ⌃H)
\begin{document}


\maketitle
\changebox

\newpage

% --- Página de resumen y modificaciones ---
\puntosclave{ %% Escribir puntos clave Apuntes CECO
    Apuntes CECO -- 33 págs.

    La distribución del tiempo sería adecuada si:
    \begin{lnum}
        \item Modelo de búsqueda de MORTENSEN y sus extensiones. Entre 4 y 6 minutos;
        \item Modelo de DIAMOND - MORTENSEN - PISSARIDES. Entre 13 y 15 minutos. Es aconsejable explicar con detalle las relaciones que describen los flujos del mercado de trabajo (función de emparejamiento, tasa de cobertura de vacantes, tasa de destrucción de empleo), así como los problemas del consumidor y la empresa. El equilibrio se puede describir en términos gráficos, subrayando la importancia de la Curva de BEVERIDGE; y,
        \item Modelos de costes de ajuste, planteando la demanda dinámica del trabajo como punto de partida (Marco Neoclásico enfoque intertemporal) y analizando las impplicaciones de la existencia de costes de ajuste. Entre 7 y 9 minutos.
    \end{lnum}
    
    }
\vspace{1cm}

\modificaciones{
     
    }

\newpage
\MyTOC
\newpage

\mylistofequations
\clearpage

\MyListOfFigures
\clearpage


\pagenumbering{arabic}
\setcounter{page}{1}


\section{Introducción}\label{intro}
La teoría del mercado de trabajo, o Economía Laboral, analiza la determinación del empleo y los salarios de equilibrio
que fijan los trabajadores (que ofertan su fuerza de trabajo) y las empresas (que demandan la
fuerza de trabajo). El mercado de trabajo se puede analizar tanto desde el punto de vista:
\begin{la}
    \item Microeconómico, estudiando el comportamiento de los agentes individuales y la fijación del
equilibrio de mercado; como,
    \item Macroeconómico, estudiando las fluctuaciones cíclicas
y tendencia del mercado de trabajo así como su influencia sobre el conjunto de la economía.
\end{la}

Mientras que los modelos empleados en el [\authorfont{Ver Tema 3.A.25}] tienen un perfil más cercano a la Microeconomía, los estudiados aquí suelen asociarse a la Macroeconomía, todo ello, con salvedades, y siendo conscientes de que no existe tal cosa corno una frontera nítida entre Macroeconomía y Microeconomía.


\subsection{Contextualización}\label{context}

\subsubsection{Relevancia}\label{importance}
El estudio del mercado de trabajo es fundamental para tratar de comprender los mecanismos de fijación del salario y el empleo de equilibrio, así como la transmisión de los posibles desequilibrios que genere.

\subsubsection{Historia}\label{history}
Debido a la gran importancia de la economía laboral, son muchos los economistas que se han centrado en su estudio:
\begin{lnum}
    \item Los autores neoclásicos aplicaron el principio de marginalidad y, en concreto, la teoría de la productividad marginal del trabajo para el estudio de la demanda de trabajo. [\authorfont{Ver Tema 3.A.25}]
    \begin{la}
        \item MARSHALL se limita a señalar algunos aspectos de la oferta y la demanda de trabajo que dejan ver que el mercado de trabajo se puede analizar como cualquier otro mercado – equilibrio parcial – tratando al mercado de trabajo separadamente. Aun así, no se encuentra una investigación sistemática en sus \textit{Principios de Economía} (1890).\footnote{%
            Para MARSHALL el desempleo no parecía ser un problema abrumador, lo que podría ser explicado por el contexto de su obra. Según MATTHEWS (1990):\\
“\textit{Unemployment, particularly in combination with inflation has made the functioning of the labor market a central topic in present-day economics. Unemployment has been judged as both intellectually anomalous and a social challenge. This emphasis is absent in Marshall. The social problem that disturbed his conscience was poverty; and poverty might have a number of causes, of which unemployment was only one.}”
        }
        \item HICKS, en su obra “\textit{Theory of Wages}” (1932), estudió la demanda de trabajo mediante la productividad marginal del trabajador. HICKS señala que la teoría económica no era capaz de explicar el desempleo, un fenómeno innegable de la realidad\footnote{HICKS escribe en el contexto de la Gran Depresión.}. Señala que de acuerdo con la teoría neoclásica los salarios deben reducirse en presencia de desempleo involuntario y llega a la conclusión de que tanto los sindicatos como el desempleo friccional juegan un rol importante\footnote{%
            HICKS en el marco del modelo IS-LL atribuye el desempleo a la rigidez salarial y deja entrever que tiene que ver con la forma de negociación salarial y el papel de los sindicatos. Esta rigidez salarial provoca un salario que no vacía el mercado de trabajo y por tanto genera paro involuntario.
        }. Sin embargo, la obra de HICKS fue eclipsada por la publicación de la Teoría General del Empleo, el Interés y el Dinero de KEYNES (1936).
    \end{la}
    \item Debido a la imposibilidad de explicar el desempleo\footnote{%
        Definición de según la Organización Internacional del Trabajo: “Las personas en desocupación, o personas desocupadas, se definen como todas aquellas personas en edad de trabajar que no estaban ocupadas, que habían llevado a cabo actividades de búsqueda de un puesto de trabajo durante un período reciente especificado, y que estaban actualmente disponibles para ocupar un puesto de trabajo en caso de que existiera la oportunidad de hacerlo.”
    } de estas teorías, se necesitaban nuevas contribuciones. Estas contribuciones realizadas entre la década de 1970 y la década de 1990 cambian los supuestos tecnológicos marshallianos y obtienen como resultado la emergencia de desempleo al permitir un ajuste no perfecto del empleo ante perturbaciones.\footnote{%
        El modelo neoclásico del mercado
de trabajo no puede explicar fenómenos como el desempleo persistente, que afecta de manera
distinta a diferentes colectivos, la existencia simultánea de trabajadores desempleados y vacantes
de puestos de trabajos, la rigidez de los salarios que no permite el vaciado del mercado o la
existencia de diferencial salarial para trabajadores similares. En efecto, existen fricciones o
imperfecciones de distinta naturaleza, que reflejan el incumplimiento de alguno de los supuestos
de partida del enfoque neoclásico. Al introducir fricciones realistas, los modelos de búsqueda
desarrollan de fom1a unificada un contexto teórico para explicar estos hechos estilizados del
mercado de trabajo.
    }[\authorfont{Ver Tema 3.A.26}]
    \item Finalmente, también era necesario ahondar en el proceso de determinación de los salarios. De manera contemporánea a los modelos anteriores, surgen modelos que explican las rigideces salariales [\authorfont{Ver Tema 3.A.27}]
\end{lnum}

\subsubsection{Problemática}\label{problem}
En el presente Tema intentaremos dar respuesta a las siguientes cuestiones:
\begin{lnum}
    \item \textit{¿Cómo estudiar el componente del desempleo que recoge a los trabajadores y/o vacantes que están transicionando de un puesto y/u ocupante a otro?}
    \item \textit{¿Cuán relevante es este desempleo friccional en el volumen total de desempleo?}
\end{lnum}


\subsection{Estructura de la exposición}\label{structure}
Como el objetivo del presente tema es centrarnos en los modelos que analizan los problemas de información y los costes de ajuste, lo que da lugar a las Teorías de la búsqueda y los costes de ajuste en la dinámica de la demanda de trabajo, dividiremos nuestra exposición en tres partes:
\begin{lnum}
    \item En primer lugar, se analizará el marco general (modelo \textit{baseline}) y las principales extensiones de las Teorías/Modelos de búsqueda, que analizan las decisiones de contratación de empleo y trabajadores en un contexto de información incompleta, teniendo en cuenta que el proceso de búsqueda de empleo y de cobertura de vacantes es costoso;
    \item En la segunda parte del tema se presenta el análisis la dinámica de la demanda de trabajo y cómo la existencia de costes de ajuste en las plantillas de las empresas afecta a la rotación laboral y puede limitar la contratación, lo que explicará por qué el empleo no se ajusta de forma instantánea.
\end{lnum}


\clearpage




\section{Modelos de búsqueda} 
Los supuesto del modelo neoclásico hacen que el mercado de trabajo alcance el equilibrio por la intersección de las curvas de demanda y oferta, y ese equilibrio sea único, estable y de pleno empleo, no existiendo espacio para el desempleo (los individuos sólo pueden estar empleados o no participar en el mercado de trabajo). No obstante, en el mercado de trabajo puede existir heterogeneidad en la mano de obra, los trabajadores y empresas se reúnen de fonna descentralizada en un costoso proceso de emparejamiento de preferencias, habilidades y necesidades, distinto para cada agente y puesto de trabajo. Este proceso, además, no es instantáneo de manera que existirán agentes que no participen en el mercado (bien porque dedican todo su tiempo al ocio o por estar desanimados, dados los costes del proceso) y agentes que participen y que estén ocupados o desempleados buscando un empleo (este último caso no tenía cabida en el análisis neoclásico). Los modelos de búsqueda fueron pioneros en la fundamentación microeconómica del desempleo y surgen en la década de los setenta para explicar las altas tasas de paro persistentes, cuestionando los supuestos neoclásicos. Tratan de analizar las consecuencias que tiene el hecho de que trabajadores y puestos de trabajo sean heterogéneos sobre el mercado de trabajo y su equilibrio. Los problemas de falta de infonnación en el mercado de trabajo cobran especial interés, como muestra el modelo de las islas de Phelps. Este modelo analiza el comportamiento de agentes que buscan empleo en el mercado de trabajo y empresas que necesitan encontrar trabajadores para sus vacantes en un contexto de información imperfecta. El supuesto de infonnación imperfecta supone que en cada mercado, ciertas perturbaciones agregadas o específicas pueden afectar a variables relevantes del mercado, como el salario nominal, de manera que cuando el agente observa un aumento de su salario nominal, éste puede ser resultado de un aumento pel nivel general de precios (perturbación a nivel agregado que no varía el valor real del salario y, por tanto, tampoco la oferta de trabajo óptima que maximiza la utilidad del agente) o un aumento específico de su salario (perturbación específica que modifica el salario real y por tanto la oferta óptima de trabajo del agente). En un contexto de info1mación imperfecta como el de nuestro agente, éste es incapaz de identificar la naturaleza de la perturbación, si es agregada o específica.

\subsection{Modelo \textit{baseline} -- MORTENSEN, 1970}
El modelo de Mortensen (1970) es uno de los pioneros,\footnote{%
    El origen fundamental de las Teorías de Búsqueda se encuentra en la aportación pionera de STIGLER, quien analizaba cómo un consumidor quería adquirir un bien homogéneo, pero había varias empresas ofreciéndolo en el mercado, cada una de ellas lanzando el producto a un precio distinto sin que el consumidor supiera con antelación a qué precio lo vendía cada una de las tiendas. La aplicación más relevante en este sentido de las teorías de búsqueda al mercado de trabajo fue realizada por Mortensen. El autor se centra en modelizar el comportamiento de un trabajador desempleado que está buscando empleo, pero que actúa en un marco de información imperfecta, pues hay una multitud de vacantes disponibles, pero cada una ofrece un salario diferente (heterogeneidad) sin que sepa de antemano cuál es el salario de cada empresa (información imperfecta).\\
Hay que destacar que estas teorías de búsqueda pueden aplicarse a multitud de ejemplos. Por ejemplo, el comportamiento de un individuo que está buscando una casa de alquiler, o está buscando la obtención de un préstamo. Sin embargo, han tenido una especial relevancia en el campo de la economía laboral. Así, estas teorías de búsqueda surgen como una alternativa a la teoría de la oferta de trabajo neoclásica, que solo recogía dos posibles tipos de agentes: empleados o inactivos.
} centrándose en determinar el salario de reserva de los trabajadores desde el lado de la oferta. Considera que  Parte de los siguientes supuestos:
\begin{lnum}
    \item El mercado de trabajo es competitivo, pero al no existir información completa, los trabajadores tienen que enfrentarse a un cierto grado de incertidumbre;
    \item Los trabajadores y puestos de trabajo son heterogéneos;
    \item La información es, además de imperfecta, costosa de obtener;
    \item Se considera como desempleados a los trabajadores que no tienen un empleo, pero que lo buscan activamente;
    \item Si el trabajador está ocupado, se supone que no busca otro empleo, es decir, no existe búsqueda desde el puesto de trabajo ( \textit{on-the-job search}).
    \item El stock de parados varía continuamente, ya que existe un flujo continuo de entrada y salida de población activa y una continua rotación en el empleo.
    \item El objetivo del trabajador será alcanzar la máxima oferta salarial.\footnote{%
        El proceso de búsqueda de la mejor oferta posible tiene unos beneficios esperados, derivados de la obtención del salario, pero también unos costes esperados de seguir buscando, asociados tanto del mismo proceso de búsqueda, como de costes de oportunidad, al no percibir un salario. El trabajador a lo largo de la búsqueda irá recibiendo una serie de ofertas salariales, y evaluará beneficios y costes esperados del proceso de búsqueda para, de entre ellas, elegir la mejor.
    }
\end{lnum}

Hay que destacar que un contexto de información imperfecta es irracional para el trabajador aceptar la primera oferta que reciba, por tanto, podemos representar mediante una función de densidad ($\phi(w)$) la distribución de los salarios para el individuo $i, w_i$:

\imagenfit{PDF-Salarios.png}{Distribución de probabilidad de la variable aleatoria $w_i\equiv$ Niveles salariales del inidividuo $i$}{Apuntes ICEX-CECO. Tema 3.A.26}

Si no existiesen costes derivados de la búsqueda de empleo, el trabajador no dejará de buscarhasta que obtenga el $W_{max}$. Sin embargo, sí existen costes.\footnote{%
    Principalmente de obtención de información, pero también costes de transporte, y un coste de oportunidad elevado, derivado del hecho de dejar de obtener un salario o de rechazar una buena oferta por seguir buscando
} Por tanto, existe un trade-off: conforme el trabajador tenga que destinar más tiempo a la búsqueda de empleo, los costes en los que incurre son mayores, pero pueden ser compensados si encuentra un empleo mejor remunerado. Así, es posible definir el salario de reserva ($W_r$) como aquél que iguala los beneficios y los costes marginales del proceso de búsqueda y a partir del cual el trabajador aceptará la oferta salarial.\footnote{%
    Para determinar el salario de reserva existen dos enfoques: el \textit{no-secuencial}, en el que el trabajador decide de antemano el número de empresas que consultará y después elegirá el salario más alto que le hayan ofrecido (STIGLER, 1961); y el enfoque secuencial, en el que el trabajador determina su salario de reserva antes de buscar empleo, de forma que aceptará las ofertas que sean al menos tan buenas como ese salario de reserva (MORTENSEN, 1970). Consideraremos este último enfoque, al ser más eficiente.
}\textsuperscript{,}\footnote{%
    La curva de $BMg$ tendrá pendiente positiva al principio porque para salarios bajos, el beneficio de seguir buscando un trabajo es mayor. Sin embargo, a partir de un determinado salario, cuanto más tiempo lleve buscando el trabajador, antes querrá encontrar un empleo y la diferencia entre el valor de búsqueda y la mejor oferta salarial disponible se reduce. El coste marginal $CMg$ tiene pendiente positiva, dado que en el caso de salarios bajos, el coste de oportunidad de seguir, buscando es bajo mientras que para salarios más altos, el coste de oportunidad de seguir buscando también se incrementa. Esto quiere que el esfuerzo de búsqueda se reduce con el salario, es decir, el valor del empleo aumenta a medida que el salario es más alto.
} Para ello, el trabajador compara los costes y los beneficios de la búsqueda, determinando el salario mínimo que está dispuesto a aceptar, lo que indicará el momento en el que dejará de buscar empleo.

\imagenfit{Eq-Busqueda1970.png}{Curvas de $BMg$ y $CMg$ asoaciadas a la búsqueda de empleo: determinación del salario de reserva $W_r$}{Apuntes ICEX-CECO. Tema 3.A.26}

Por tanto, formalmente el trabajador maximiza su utilidad esperada de buscar trabajo, representada por una función de tipo vN-M, sujeto a la distribución salarial existente:

\eqblock{
\begin{aligned}
    &\max_{u}{E[u(w)]} \qquad \text{s.a.}\quad \phi(w)\\[2em]
    \Longrightarrow \text{Solución:}\quad &W_r^*=f(CMg, BMg, \alpha,t,e,\sigma), \quad \text{\scriptsize Donde,} 
    \begin{cases}
        CMg\;\;y\;\;BMg&\equiv \text{\scriptsize Curvas del agente $i$}\\
        \alpha&\equiv \text{\scriptsize Tasa de descuento de los agentes}\\
        t&\equiv \text{\scriptsize Tiempo desempleado/en búsqueda}\\
        e&\equiv \text{\scriptsize Edad del trabajador $i$}\\
        \sigma&\equiv \text{\scriptsize Cuantía del subsidio por desempleo}\\
    \end{cases}\\[1em]
    \text{Y,} \quad &P[W_r\le w]=1-P[w<W_r]\quad \sim\quad  1-\int_{W_{min}}^{W_{r}}\phi(w)\mathrm{d}w=\int_{W_r}^{W_{max}}\phi(w)\mathrm{d}w
\end{aligned}
}{Modelos de búsqueda: Modelo \textit{baseline}, MORTENSEN}

Esta área es la probabilidad de recibir una oferta salarial mayor o igual que el salario de reserva, tal que:
\imagenfit{CDF-W_r.png}{Probabilidad de $W_r\le w$}{Apuntes ICEX-CECO. Tema 3.A.26}

Como vemos, existen muchos factores que inciden en el valor de $W_r$, de forma que:\footnote{%
    Existen otros factores que pueden influir en el salario de reserva, como la tasa a la que los desempleados encuentran un trabajo, que refleja la dificultad de la búsqueda y resume las fricciones del mercado, las características del trabajador (como su cualificación) y el esfuerzo realizado en el proceso de búsqueda. Cuanto mayor sea la tasa, mayor será el salario de reserva. Este factor se analizará más adelante.
}
\begin{la}
\item Afectará positivamente al salario de reserva $W_r$:
    \begin{lnum}
        \item Un aumento de los beneficios marginales de seguir buscando provocará cambios en el salario de reserva y en las opciones de búsqueda\footnote{%
            Esto se debe a que los salarios reales son procíclicos, por lo que los beneficios de seguir buscando empleo aumentan en las expansiones y disminuyen en las recesiones.
        } y, por tanto, cuando la curva BMg se desplace a la derecha, tanto el salario de reserva como la duración de la búsquedaaumentarán, siendo $\partial W^*/\partial BMg>0$.
        \item El subsidio por desempleo (el cual puede aproximarse al coste de oportunidad de la búsqueda) afecta positivamente al salario de reserva porque reduce los costes marginales de seguir buscando, siendo $\partial W^*/\partial \sigma>0$.
    \end{lnum}
\item Por otro lado, afectará negativamente al salario de reserva $W_r$:
    \begin{lnum}
        \item Un aumento en los costes marginales de seguir buscando (desplazamiento de la curva CMg hacia arriba) provocará una disminución del salario de reserva, lo cual reducirá la duración de la búsqueda y, por tanto, el desempleo, siendo $\partial W^*/\partial CMg<0$.
        \item Si los trabajadores tienen tasas de descuento elevadas se producirá un desplazamiento hacia abajo de la curva BMg, dado que el trabajador tiene preferencia por el presente y los beneficios de seguir buscando son menores. De esta forma, el salario de reserva y la duración de la búsqueda serán menores, siendo $\partial W^*/\partial \alpha<0$.
        \item Cuanto mayor sea el tiempo de búsqueda, menores serán los beneficios, debido a la pérdída de cualificación y desmotivación, y menor será el salario de reserva, por lo que: $\partial W^*/\partial t<0$.
        \item Por motivos similares al anterior factor, el salario de reserva depende inversamente de la edad del trabajador, tal que $\partial W^*/\partial e<0$.
    \end{lnum}
\end{la} 

De todas ellas, la que más atención acapara en la literatura es el seguro por desempleo $\sigma$. Al reducir los costes marginales de búsqueda e incrementa el salario de reserva, aumenta también la duración del paro friccional.\footnote{%
    Por ejemplo, si recibe una oferta de 500 euros y no hay seguro de desempleo, renunciar a ella implica renunciar a 500 euros; pero si tengo un seguro de desempleo de 200 euros, entonces renunciar a la oferta solo me supone una pérdida de 300 euros
} ya que, como seguir en paro es menos costoso, puedo esperar mucho hasta ponerme a trabajar. Por tanto, nos permite observar que la configuración institucional importa en el funcionamiento del mercado de trabajo y que el seguro de desempleo tiene un claro componente de riesgo moral, al dar incentivos a los trabajadores a no ocupar un empleo. Al mismo tiempo, el seguro de desempleo tiene efectos positivos muy importantes sobre el bienestar de los parados y sus familias que logran una mayor estabilidad de sus ingresos (agentes aversos al riesgo) y \footnote{%
    Al mismo tiempo, a nivel macroeconómico, este seguro por desempleo es un estabilizador automático importante que suaviza las fluctuaciones cíclicas. Por último, es positivo por motivos de eficiencia, pues permite a los trabajadores utilizar el tiempo necesario para encontrar un empleo que se ajuste a sus preferencias y habilidades.
}  evita que los trabajadores deban que precipitarse. Por ello, al tener el seguro por desempleo importantes beneficios, debemos centrarnos en mitigar sus costes, en particular:
\begin{lnum}
    \item Cobertura decreciente en el tiempo, para que progresivamente el CMG de búsqueda vaya siendo mayor y se reduzca el salario de reserva
    \item Condicionalidad: los receptores están obligados a participar en cursos, talleres, entrevistas, procesos selectivos, etc.
    \item Imposibilidad de rechazar ofertas de empleo
\end{lnum}

\subsubsection{Implicaciones de política económica}
Este modelo de búsqueda permite extraer algunas conclusiones relevantes sobre el funcionamiento del mercado de trabajo y la existencia de desempleo. Se observa que cuando los trabajadores perciben correctamente la distribución de salarios, su desempleo es eficiente en el sentido de que los trabajadores siguen estrategias óptimas y optan voluntariamente por permanecer desempleados, rechazando ofertas salariales menores que su salario de reserva. Este modelo explica, por tanto, el desempleo voluntario de la teoría neoclásica. Sin embargo, en la realidad también existe paro involuntario y persistente, que este modelo no es capaz de explicar. Por ello han surgido extensiones que pretenden dar respuesta a esta evidencia empírica.

\subsection{Modelo de referencia: flujos y emparejamiento}
El modelo de búsqueda visto en el apartado anterior se ha centrado únicamente en el lado de la oferta por lo que no se considera ni la demanda de trabajo, ni los determinantes salariales. Sin embargo, las teorías modernas de búsqueda analizan tanto el lado de la oferta como el de la demanda. En efecto, por el lado de la demanda, la existencia de los puestos de trabajo vacantes es un elemento fundamental para analizar el emparejamiento que se produce en el mercado, entre demandantes, vacantes e ineficiencias. Esto nos permitirá introducir la existencia de desempleo involuntario. 

PISSARIDES\footnote{La base teórica de este modelo se fundamenta en PISSARIDES, 2000} crtica la Teoría del salario de reserva porque considera que realmente los trabajadores no inciden sobre su oferta salarial sino que deciden sobre las vacantes existentes en el mercado. Al considerar al mismo tiempo el lado de la oferta y de la demanda, este modelo permite analizar en profundidad los flujos de desempleo y la interacción entre los agentes. Al introducir la demanda de trabajo en un contexto de heterogeneidad de trabajadores y puestos de trabajose considera que las empresas, paralelamente, establecen un criterio de reserva para la selección de candidatos. De esta forma, sólo contrarán al candidato que supere o iguale este criterio de reserva.

Esto ha dado lugar a un modelo conocido como el modelo Diamond-Mortensen-Pissarides (DMP)\footnote{%
    Precisamente estos tres economistas recibieron en 2010 el Premio Nobel de Economía \textit{Por su análisis de los mercados con fricciones de búsqueda..}
}, que se ha convertido en el modelo de referencia para analizar el desempleo, la formación de los salarios y la existencia de puestos de trabajo vacantes. Según estos autores, la duración del paro depende de la intensidad de la búsqueda y de los costes de ajuste. Se trata de un modelo de eficiencia en la búsqueda que tiene en cuenta la creación y la destrucción de empleo y analiza el emparejamiento entre trabajadores y vacantes.

\subsubsection{Supuestos}
Para su análisis supondremos que:
\begin{lnum}
    \item Tanto empresas como trabajadores tienen pleno conocimiento de los procesos emparejamiento pero no tienen interés en coordinar sus acciones;
    \item Hay muchas empresas de pequeño tamaño, de forma que cada empresa podrá ofrecer sólo una vacante, y muchos trabajadores, $L$, con vida infinita, en un modelo de competencia atomizada. 
    \item El equilibrio agregado del modelo será aquél que garantiza que tanto empresas como trabajadores maximicen sus funciones objetivo, sujeto a la tecnología de emparejamiento, y para el cual el flujo de trabajadores que encuentran un empleo es igual al mismo flujo de trabajadores que lo pierden;
    \item Además, sólo los trabajadores parados buscan empleo.
\end{lnum} 

Este modelo se basa por tanto en la existencia de importantes fricciones que afectarán tanto a los trabajadores como a las empresas; de hecho, recoge todas las fricciones que puedan afectar al mercado de trabajo, no solo son informativas. Éstas llevan a que se tarde tiempo en emparejar oferta con demanda (los trabajadores enviarán currículums, preguntarán a conocidos, irán a entrevistas, etc. las empresas anunciarán la vacante, llamarán a candidatos, les entrevistarán, etc.), por lo que puede haber perfectamente trabajadores desempleados y vacantes sin cubrir, dando lugar a desempleo friccional. Sin embargo, parece evidente que cuantos más desempleados haya y más vacantes existan mayores serán los emparejamientos que se producirán en la economía (tercer supuesto).\footnote{%
    Esta es la idea de la función de matching o de emparejamiento, que guarda una evidente similitud con la función de producción, los inputs serán las vacantes y los desempleados, el output será el número de emparejamientos o de colocaciones, y en lugar de depender de la tecnología dependerá de las fricciones que haya en el mercado de trabajo.
} 

\subsubsection{Función de emparejamiento} 
Definiendo $u$ como la tasa de desempleo (fracción de trabajadores desempleados) y $v$ como la tasa de vacantes (número de puestos vacantes entre la población activa), asumimos que $uL$ es el número de trabajadores desempleados y $vL$ es el número de vacantes para emparejar.\footnote{%
    Asumimos además que la función sigue rendimientos constantes de escala, que es una propiedad importante, tal y como se ha demostrado empíricamente.
} 

De esta forma, podemos definir la función de emparejamiento, como la que relaciona las vacantes y la tasa de paro, de forma que:

\eqblock{
\begin{aligned}
    &\text{Función de emparejamiento:} \quad \boxed{m(u,v)=m},\qquad \text{\scriptsize donde}
    \begin{cases}
        m\in(0,1)&\equiv\text{\scriptsize Función/Tasa de emparejamiento}\\
        u\in(0,1)&\equiv\text{\scriptsize Tasa de desempleo}\\
        v\in(0,1)&\equiv\text{\scriptsize Tasa de vacantes}\\
        L&\equiv\text{\scriptsize Población activa}\\
    \end{cases}\\[1em]
    &\text{asumimos que: }\;m(u,v)\;\text{ es }\underbrace{\text{Homogenea de grado $n=1$}}_{\text{\tiny (Recordar, $\equiv F(\lambda x_i)=\lambda^nF(x_i)$)}}\quad \Longrightarrow \quad Lm(u,v)=m(uL,vL)\\[1.5em]
    &\text{Y,} \qquad \boxed{\partial m/\partial u >0\quad;\quad \partial m/\partial v<0}\\[1em]
\end{aligned}
}{Desarrollo análitico. Función de Emparejamiento y derivadas.}

\eqblock{
\begin{aligned}
    &\text{(1)}\quad \text{Definimos la } \text{\textbf{Tasa de cobertura de vacante} ($q\in(0,1)$ como:}\\[0.5em]
    &q=\frac{m(uL,vL)}{vL}\underbrace{=}_{\text{\tiny Por $H(1)$}} \frac{m(u,v)}{v} \underbrace{\to}_{\text{\tiny Reescribiendo, con $\begin{cases}u=v\;\cdot\;(u/v)\\v=v\;\cdot\;1\end{cases}$}}q=\frac{m(v\;\cdot\;(u/v),\;v\;\cdot\;1)}{v}\underbrace{\to}_{\text{\tiny Utilizando $\lambda =v$}}q=m\left(\frac{u}{v},1\right)\\[1em]
\end{aligned}
}{Función de cobertura de vacante}

De esta forma, podemos ver que la intuición económica persiste en el desarrollo analítico:
\begin{la}
    \item Si la proporción entre desempleados y vacantes es alta ($\uparrow \frac{u}{v}$), la proporción entre vacantes y desempleados es baja ($\downarrow \frac{v}{u}$), será más difícil para las empresas encontrar trabajadores, y por tanto la tasa de cobertura de vacantes $q$ será baja.
    \item Si la proporción entre desempleados y vacantes es baja ($\downarrow \frac{u}{v}$), la proporción entre desempleados y vacantes es alta ($\uparrow \frac{v}{u}$), las vacantes serán escasas y $q(\theta)$ será alto, puesto que el emparejamiento será más fácil.
\end{la}

\eqblock{
\begin{aligned}
    &\text{(2)}\quad \text{Definimos la } \text{\textbf{Probabilidad de encontrar un empleo} ($f\in(0,1)$ como:}\\[0.5em]
    &f=\frac{m(uL,vL)}{uL}\underbrace{=}_{\text{\tiny Por $H(1)$}} \frac{m(u,v)}{u} \underbrace{\to}_{\text{\tiny Reescribiendo, con $\begin{cases}u=u\;\cdot\;1\\v=u\;\cdot\;(v/u)\end{cases}$}}
    f=\frac{m(u\;\cdot\;1,\;u\;\cdot\;(v/u))}{v}\underbrace{\to}_{\text{\tiny Utilizando $\lambda =u$}}f=m\left(1,\frac{v}{u}\right)
\end{aligned}
}{Función de probabilidad de encontrar un empleo}

Por tanto, de forma similar al análisis anterior, los desempleados encontrarán trabajadores más fácilmente cuando haya más puestos de trabajo relativo al número de trabajadores disponibles:
\begin{la}
    \item Si la proporción entre vacantes y desempleados es alta ($\uparrow \frac{v}{u}$), a un desempleado le resultará más fácil obtener un empleo ($f$ será alto).
    \item Si la proporción entre vacantes y desempleados es baja ($\downarrow \frac{v}{u}$), a un desempleado le resultará más dificil obtener un empleo ($f$ será alto).
\end{la}

Definiendo $\theta$ como el grado de presión en el mercado:
\eqblock{
\begin{aligned}
    &\left(\text{$\theta=\frac{v}{u}\equiv$ \textbf{Presión del mercado de trabajo}, tenemos:}\right)\\[1em]
    &\hspace{2em}\begin{cases}
        &\text{(1)}\qquad \boxed{q(\theta)=m\left(\frac{1}{\theta}, 1\right)}\quad \text{\scriptsize Donde,}\quad 
    \begin{cases}
        \partial q/\partial \theta\leq0; \quad \text{y, por tanto,}\\
        \epsilon=\frac{\partial q}{\partial \theta}\cdot\frac{\theta}{q}\in (-1,0]
    \end{cases}\\[3em]
    &\text{(2)}\qquad \boxed{f(\theta)=m\left(1, \theta\right)}\quad \text{\scriptsize Donde,}\quad 
    \begin{cases}
        \partial f/\partial \theta>0; \quad \text{y, por tanto,}\\
        \epsilon=\frac{\partial f}{\partial \theta}\cdot\frac{\theta}{f}\in (0,1)
    \end{cases}
    \end{cases}\\[1em]
    &\text{Y, por $H(1)$:}\\[0.5em]
    &\hspace{2em}\boxed{q(\theta)=m\left(\frac{1}{\theta}, 1\right)\underbrace{\to}_{\text{\scriptsize Multiplicamos por $\theta$}}\theta q(\theta)=m\left(1, \theta\right)=f(\theta)}
\end{aligned}
}{Relación opuesta subyacente entre $f(\theta)$ y $q(\theta)$. La presión en el mercado de trabajo.}

La dependencia de estas funciones del número de personas del mercado ($L$) y su grado de rigidez o presión ($\theta$) supone un elemento crucial en el análisis:

$$\begin{aligned}
    &q(\theta)=\frac{f(\theta)}{\theta} \underbrace{\Longrightarrow}_{\substack{\text{\scriptsize Para preservar las condiciones} \\\text{\scriptsize de derivación: $q'(\theta)<0$}}}f'(\theta)-f(\theta)<0 \iff \frac{f'(\theta)\theta}{f(\theta)}<1\\[1em]
    \text{Por tanto:}
    &\begin{cases}
        \epsilon_{f,\theta}\in(0,1)\quad \text{y}\\
        \epsilon_{q, \theta}\in(-1,0]
    \end{cases}
\end{aligned}$$

dando lugar a una externalidad que proviene del hecho de que durante el proceso de emparejamiento, existe una probabilidad de que la empresa no encuentre un trabajador y de que el desempleado no encuentre un trabajo durante un momento de tiempo, sea cual sea el salario. 

$\Longrightarrow$ El salario no es el único mecanismo de ajuste, sino que el número de personas del mercado también puede influir en esa probabilidad. Este hecho se conoce con el nombre de externalidad de búsqueda o congestión.

\paragraph{Equilibrio estacionario}
Teniendo en cuenta las ecuaciones planteadas, debemos introducir un último parámetro para encontrar la condición de equilibrio estacionario\footnote{%
    El modelo original parte de un enfoque dinámico, en el que se tiene en cuenta la evolución de estos flujos a lo largo del tiempo. Sin embargo, por simplicidad analítico nos centraremos en el caso estacionario, de modo que la evolución del desempleo sea $0$.
}. 

Existe una probabilidad exógena de que un trabajador pase del empleo al desempleo definida como $\lambda$, que es equivalente a la tasa de destrucción de empleo.\footnote{%
    Este proceso de destrucción es independiente del proceso por el cual un desempleado se incorpora a una vacante. Sin cambios en la población activa, el número medio de trabajadores que pasa a estar desempleado en un periodo de tiempo será $\lambda(1 - u)L\;\mathrm{d}t$ y el número medio de los que encuentran un empleo es $mL\; \mathrm{d}t$. El número medio de trabajadores que encuentran un empleo también se puede escribir como la probabilidad de encontrar un empleo por el número de desempleados: $uLf(\theta)=uL\theta q(\theta)$ donde $u\theta q(\theta) \;\mathrm{d}t$ es la probabilidad de transición de un desempleado.
} Por lo tanto, la evolución del desempleo medio viene dada por la diferencia de ambos flujos, empleados que entran al desempleo y desempleados que salen hacia el empleo, siendo el estado estacionario ($\dot u$) áquel que cumple:

\eqblock{
\begin{aligned}
    \text{Evolución del desempleo:} \quad &\dot u=\underbrace{\lambda(1-u)}_{\text{\scriptsize Flujo de entrada: FED}}-\underbrace{\theta q(\theta)u}_{\text{\scriptsize Flujo de salida: FSD}} \\[1em]
    \text{Y, Estado Estacionario en:}\quad &\dot u=0\iff\underbrace{\lambda(1-u)=\theta q(\theta)u}_{\text{FED = FSD}}\\[2em]
    \text{Entonces,}\Longrightarrow\text{$u$ de equilibrio es:}\quad &u(\lambda, \theta)=\frac{\lambda}{\lambda+\theta q(\theta)}\quad \text{y}\quad \forall(\lambda, \theta), \quad \exists!u
\end{aligned}
}{Equilibrio Estacionario: Tasa de desempleo de equilibrio}

Precisamente por ello, podemos obtener la relación inversa\footnote{%
    ¿Cómo se observa esa relación inversa? Por la probabilidad de que un desempleado encuentre trabajo: a mayor número de vacantes, mayor probabilidad y por tanto menos desempleados y viceversa.
} entre desempleo ($u$) y vacantes ($v$) dado una tasa de destrucción de empleo (exógena) determinada ($\bar\lambda$):\footnote{%
    En honor a WILLIAM BEVERIDGE, quien había descubierto la relación inversa entre vacantes y desempleados a principios del siglo XX.
}

\imagenfit{Curva-Beveridge.png}{Curva de Beveridge. Fijando: $\bar \lambda$}{Apuntes ICEX-CECO. Tema 3.A.26}

Esta curva de Beveridge nos permite conocer cómo de ineficiente es el mercado de trabajo de una economía por la existencia de esas fricciones que impiden el emparejamiento entre vacantes y desempleados. Más concretamente, para un mismo nivel de vacantes, el nivel de fricción determinará la distancia de la Curva al origen: cuanto más alejada se encuentre del origen más fricciones existen en el mercado, mientras que cuanto más cercana al origen esté,  menos habrá.\footnote{%
    Démonos cuenta que en un entorno neoclásico no habría ni desempleados buscando trabajo ni vacantes sin ocupar, por lo que realmente no habría Curva de BEVERIDGE, pero como las fricciones son una realidad, la Curva de BEVERIDGE es una representación de estas fricciones.
} 

Sin embargo, aunque esta relación puede extraerse de las ecuaciones presentadas, la obtención cocnreta del equilibrio requiere del cumplimiento y desarrollo de dos ecuaciones adicionales.

\paragraph{Condición del trabajador: salario de reserva}
Los trabajadores generalmente afectan al equilibrio a través de su búsqueda de trabajo y su influencia en la determinación del salario. En este caso, se pretende derivar el rendimiento de un trabajador desempleado y empleado. Un trabajador gana un salario $w$ cuando está empleado y busca un puesto cuando está desempleado, obteniendo un rendimiento real $z$ constante e independiente de los rendimientos del mercado, que obtiene cuando está buscando un puesto (por ejemplo, por medio de prestaciones por desempleo o el rendimiento de las actividades de ocio). $U$ y $W$ son el valor descontado al presente de la renta esperada de un desempleado y un empleado, respectivamente. Por lo tanto, podemos establecer las siguientes relaciones:

\eqblock{
\begin{aligned}
    \underbrace{\text{Salario reserva:}}_{\substack{\text{\tiny Se trata del rendimiento esperado del capital humano}\\\text{\tiny del buscador de trabajo mientras está desempleado}}} \quad &rU=z+\theta q(\theta)(W-U)\qquad \equiv\text{\scriptsize Compensación min necesaria para abandonar desempleo}\\
    \underbrace{\text{Renta salarial:}}_{\substack{\text{\tiny Renta percibida por el empleado teniendo en cuenta}\\\text{\tiny el salario $w$ y renta adicional dado el riesgo de desempleo}}} \quad &rW=w+\lambda(U-W)\qquad \equiv\text{\scriptsize Compensación min necesaria para abandonar el empleo}\\
    \text{Entonces, si} \quad &\not\exists \;\textit{on-the-job search} \Longrightarrow \text{trabajador empleado $\iff \underbrace{W\ge U \Longrightarrow w\ge z}_{\text{\scriptsize Salario superior a rendimiento desempleo}}$}
\end{aligned}
}{Condición del trabajador (salarial): $W\ge U$}

\paragraph{Condición de creación de empleo: Producto}
Nos indica cuántas vacantes serán creadas en la economía por parte de las empresas. La creación de una vacante tiene un beneficio esperado determinado, que será el beneficio si la vacante se cubre y un trabajador pasa a producir. Sin embargo, la misma vacante tiene un coste si se queda vacía, pues habrá que incurrir en un proceso costoso para cubrirla. 
Esta condición nos dice entonces que la empresa creará vacantes hasta el punto de que el beneficio esperado de la última vacante iguale al coste de creación de la vacante. 

\eqblock{
\begin{aligned}
    \underbrace{\text{Coste de búsqueda:}}_{\substack{\text{\tiny Es el coste de un puesto no cubierto $pc$ menos el retorno}\\\text{\tiny neto del trabajador por la probabilidad de curbir vacante}}} \quad &rV=-pc+q(\theta)(J-V)\quad \to\hspace{0.2em}
    \begin{aligned}
        &\substack{\text{\tiny En equilibrio: $V=0$, obtiene todos}\\\text{\tiny B$^o$'s de cubrir vacantes}}\\
        &\text{\tiny Y, $\quad J=\frac{pc}{q(\theta)}$}\\
        &\text{\tiny Por tanto, duración esperada de una vacante: $\frac{1}{q(\theta)}$}
    \end{aligned}\\[1em]
    \underbrace{\text{Valor empleado:}}_{\substack{\text{\tiny Se trata del coste neto del puesto debe ser igual a el coste bruto, $p-w$ }\\\text{\tiny menos la probabilidad de perder el empleado por el valor del puesto ocupado $J$}}} \quad &rJ=p-w-\lambda J\\[1em]
    \Longrightarrow\text{Condición de creación de empleo:} \qquad &\boxed{p-w-\frac{(r+\lambda)pc}{q(\theta)}=0}
\end{aligned}
}{Condición de creación empleo}

\paragraph{Equilibrio negociador: Fijación del salario}
Cuando una empresa y un trabajador entran en contacto se abre un proceso de negociación importante.\footnote{%
    Para ambos agentes aceptar el matching tiene unos beneficios obvios: para el trabajador la posibilidad de lograr un empleo que le aporte unos ingresos salariales, y para la empresa la posibilidad de tener una vacante ocupada que le genere unos ingresos.\\
Al mismo tiempo, rechazar el matching tiene unos costes: para el trabajador el coste de tener que seguir desempleado y continuar buscando, y para la empresa el coste de tener una vacante sin ocupar y tener que seguir llamando a candidatos para realizar entrevistas.
}
Es evidente que ambos agentes tienen incentivos a llevar a cabo el emparejamiento, pues están mejor haciéndolo que rechazándolo. Pero no por ello dejarán de  negociar cuál será la contraprestación del contrato: \textit{el salario}.\footnote{%
    Principal diferencia con el enfoque previamente visto, Teoría de búsqueda, donde los salarios venían dados.
}
Éste debe ser lo suficientemente elevado como para que el trabajador quiera aceptar el empleo, y lo suficientemente reducido como para que la empresa quiera aceptar al trabajador; pero, además, como el emparejamiento es beneficioso hay un margen de salarios que satisfacen ambas condiciones.\footnote{%
    Una implicación importante de este modelo es que, dado que existen costes de búsqueda para la empresa y el trabajador, el emparejamiento y la ocupación de los puestos de trabajo existentes produce rentas. Esto supone que en el caso en que la relación entre empleador y trabajador se rompa, uno de los dos agentes (y a menudo ambos) empeoran su utilidad. Esas rentas dan a los empleadores cierto poder de mercado frente a los trabajadores, lo que implica que, si se produce un pequeño recorte salarial, no todos los trabajadores dejarán sus puestos de trabajo dado que es difícil encontrar un nuevo puesto de trabajo. Un corolario de esto es que hay trabajadores de la misma calidad que terminan siendo pagados con diferente salario si están empleados en distintas empresas.
}

Como consecuencia de ello, en el equilibrio $V=0$, se deriva el salario tanto de la condición de creación de trabajo
como de la condición del trabajador, acorde con una solución de negociación de Nash para la que el salario maximiza el producto ponderado del rendimiento neto del trabajador y la empresa en el
proceso de emparejamiento.\footnote{%
    Cada agente tratará de obtener el mejor resultado para sí mismo, sabiendo que si no llegan a un acuerdo tienen que volver a pasar por un costoso proceso de búsqueda, lo que le permite amenazar al otro agente.
}. Por tanto, el salario debería componerse de un salario de reserva ($rU$) y una parte que refleje la ganancia neta que se crea al aceptar un puesto de trabajo ($\beta(p-rU)$), para lo que se introduce el parámetro p, que mide la fuerza de negociación del trabajador, esto es, la participación del trabajo en el beneficio total que obtiene la empresa de un puesto ocupado.
\eqblock{
\begin{aligned}
    &w=\underbrace{rU}_{\text{\scriptsize Salario de reserva}}+\underbrace{\beta(p-rU)}_{\text{\scriptsize Ganancia neta que se crea al aceptar}}\qquad \text{[1]}\\[1em]
    &\text{Sustituyendo en [1], la ecuación del salario de reserva,} \quad (\text{\scriptsize $rU=z+\theta q(\theta)(W-U)$})\\[1em]
    &\text{Ecuación de salarios:}\qquad \underbrace{\boxed{w=(1-\beta)z+\beta p(1+c\theta)}}_{\text{\scriptsize $\sim$ Curva de Oferta del Trabajo. Marco Neoclásico}}
\end{aligned}
}{Ecuación de salarios}

El equilibrio en términos de salario y grado de rigidez del mercado de trabajo se puede hallar a través de la ecuación de salarios y la condición de creación de empleo. La ecuación de salarios implica que a mayor grado de rigidez del mercado de trabajo, mayor salario (los trabajadores tendrán más poder de negociación). Por su parte, cuanto mayor sea el salario, mayor será el beneficio para la empresa de abrir una vacante, por lo que la condición de creación tiene pendiente negativa en el espacio de vacantes y grado de estrechez.

\imagenfit{Ecuacion_Salarios.png}{Curvas de condición salarial y ecuación de salarios}{Apuntes ICEX-CECO. Tema 3.A.25}

\subsubsection{Implicaciones de política económica}

Una vez que ya hemos configurado nuestro modelo de matching y hemos logrado alcanzar el equilibrio en el mismo podemos preguntarnos, ¿qué implicaciones de política pública supone este modelo?

\paragraph{Fricciones y velocidad de ajuste}
Los modelos de búsqueda han sido utilizados también para analizar cómo las perturbaciones agregadas son transmitidas en el mercado de trabajo a través de la creación y destrucción de empleo y las fluctuaciones cíclicas del desempleo. Las fricciones en la búsqueda pueden predecir por qué ante una perturbación negativa el incremento en el desempleo se produce más rápido que la caída del desempleo cuando se da un shock positivo en la economía. Esto se debe a que ante un shock adverso se produce un incremento inmediato en los despidos y un aumento en el desempleo, mientras que un shock positivo solo genera una caída gradual del desempleo porque la contratación de empleados es costosa en términos de tiempo.
 
Por otro lado, como ya hemos señalado, la curva de BEVERIDGE relaciona las tasas de paro y las vacantes (medidas en términos de la población activa) cuando el desempleo pennanece estable, es decir, cuando los flujos de entrada y salida del mismo son idénticos:\footnote{%
    Esta curva se obtiene directamente de la ecuación de paro de equilibrio, para la cual el modelo DMP proporciona una explicación teórica y la ubicación de la economía en la curva. Tiene pendiente negativa porque al disminuir las vacantes, la probabilidad de encontrar trabajo ($f$) disminuye. La curva se desplazará hacia la derecha siempre que $\lambda$ aumente.\\Los puntos de la recta de 45° representan combinaciones en las que el número de parados es igual al número de vacantes. El paro asociado al nivel $u_0$ sería un paro friccional, resultado de un ajuste incorrecto del mercado, que puede ser interpretado como el correspondiente a la tasa natural de paro si los salarios son los de equilibrio. Los puntos situados a la derecha de la recta de 45° supondrían un número de parados mayor al de vacantes, por lo que existiría un exceso de oferta de trabajo o unos salarios rígidos y altos. Por su parte, los puntos situados a la izquierda de la recta supondrían un número de vacantes superior al de desempleados. El paso de un punto como $X$ a $Y$ supondría, en este contexto, un aumento del paro.
}

\imagenfit{Curva-Beveridge.png}{Curva de Beveridge. Fijando: $\bar \lambda$}{Apuntes ICEX-CECO. Tema 3.A.26}

La cuestión fundamental es determinar qué factores desplazan la curva de Beveridge:\footnote{%
    Según BLANCHARD y DIAMOND (1989) existen tres clases de perturbaciones que afectan a la posición de la curva: shocks de actividad agregada (que produce movimientos a lo largo de la curva al mover el desempleo y las vacantes en dirección contraria), shocks de reasignación (que afectan en el mismo sentido al paro y a las vacantes y producen desplazamientos de la curva) y shocks de la oferta de trabajo (que pueden aumentar la efectividad del emparejamiento al aumentar el paro y disminuir las vacantes).
} 
\begin{la}
    \item La curva se desplaza hacia la derecha en respuesta a todos aquellos factores que causen una reducción de la eficiencia en la búsqueda o a la existencia de rigideces en el mercado de trabajo que empeoren el funcionamiento del mercado y dificulten el ajuste.\footnote{%
        Un ejemplo clásico es el desajuste o mismatch en el emparejamiento. Esto puedo producirse como consecuencia de una distribución de salarios desigual entre trabajadores de una economía, a la posible falta de movilidad sectorial o regional de los trabajadores o a la falta de adecuación entre la cualificación de los trabajadores y los requisitos de los puestos. En general, este tipo de perturbaciones provocan desplazamientos hacia fuera de la curva, porque las tasas de equilibrio de vacantes y desempleo aumentan simultáneamente.
    }; y,
    \item Por el contrario, la Curva de BEVERIDGE se desplazará hacia la izquierda cuando se incremente la eficiencia del emparejamiento en el mercado.
\end{la}

En términos del modelo, la condición de creación de empleo se puede expresar como una función lineal entre vacantes y desempleo, ya que, como se deduce del gráfico anterior, el grado de rigidez del mercado de trabajo es independiente del desempleo. De este modo, la interacción entre la curva de Beveridge y la condición de creación de empleo permitiría obtener las vacantes y el desempleo de equilibrio:
\begin{la}
    \item La curva de creación de empleo recoge los factores de demanda de trabajo: cuanto mayor sea el número de desempleados, dado un grado de rigidez, mayor será el número de vacantes que se abren.
    \item La curva de Beveridge comprende los factores de demanda: cuando las contrataciones son muy elevadas el número de desempleados será menor.
\end{la}

\imagenfit{Vacantes-Desempleo_Equilibrio.png}{Condición de creación de empleo como función lineal de $v,u$}{Apuntes ICEX-CECO. Tema 3.A.26}

El equilibrio dependerá de los parámetros del modelo.\footnote{%
    La productividad del trabajo, el coste de las vacantes, el nivel por prestaciones de desempleo y el poder de negociación sólo afectan a la curva de creación de empleo. Por su parte, los parámetros que determinan el emparejamiento (elasticidad y eficiencia del mismo) y la destrucción de empleo afectan a las dos curvas.
} 

Tradicionalmente, la literatura ha señalado dos principales tipos de políticas intervencionistas en el mercado de trabajo:
\begin{lnum}
    \item \textit{Políticas activas.} Orientadas a incrementar la probabilidad de que los trabajadores desempleados encuentren un empleo cuanto antes.
    \item \textit{Políticas pasivas.} orientadas a proteger al desempleado durante la duración del desempleo para evitar que se genere una caída drástica de sus ingresos y de su bienestar.
\end{lnum}

Este modelo de matching nos permite analizar muy fácilmente las implicaciones que este tipo de políticas tendrían sobre el mercado de trabajo.

\paragraph{Políticas pasivas}
En el modelo DMP, unas prestaciones por desempleo mayores se asocian a mayor desempleo, no porque el trabajador sea más perezoso, sino porque su posición en la negociación es mejor, dejando un menor margen al empleador, que reduce el número de vacantes.\footnote{%
    Los aumentos del poder de negociación de los trabajadores, de la protección por desempleo y de los costes de creación de vacantes tienen el efecto contrario y desplazan la curva de creación de empleo hacia la derecha, alcanzándose un nuevo equilibrio con menos vacantes y mayor desempleo. Esto se debe a que se reduce el beneficio de abrir una vacante para la empresa, ya que se incrementa el coste del trabajador (aumenta el salario al elevarse el poder de negociación o las prestaciones por desempleo) o directamente al incrementarse el coste de la vacante.
} 

Esto genera una presión al alza en los salarios y una presión a la baja en las vacantes, que es lo que se conoce como \textit{efecto macro}.\footnote{%
    \authorfont{HAGEDOM et al.} (2016) demuestran empíricamente que este efecto macro ocurrió en Estados Unidos tras la crisis financiera, donde la duración de la prestación se incrementó. A través de la comparación de regiones limítrofes, demostraron que unas prestacionesmás generosas en un estado generaban un desempleo mayor vía salarios más altos y menores vacantes.
} 

\paragraph{Políticas activas}
Hasta ahora nos hemos movido a lo largo de la curva de Beveridge (lo que suele estar ligado al ciclo económico), pero puede ser más deseable mover la curva de Beveridge (lo que está ligado a la eficiencia de emparejamiento y, por tanto, la eliminación de fricciones). Más concretamente, un desplazamiento hacia el origen de esta curva vendría a reflejar un mejor funcionamiento del mercado de trabajo al reducir las fricciones existentes en el mismo, lo que permite emparejar mejor vacantes y desempleados.

Un incremento de la productividad, aumenta el incentivo de la empresa a sacar vacantes, de modo que para cada nivel de desempleo hay más vacantes. La curva de creación de empleo se desplaza a la izquierda hasta alcanzar un nuevo equilibrio con menor desempleo y mayor número de vacantes. 

En el caso de un aumento de la tasa de destrucción de empleo:
\begin{lnum}
    \item La curva de Beveridge se desplaza hacia fuera, ya que aumentan los flujos de entrada en el desempleo en relación con los de salida.
    \item La curva de creación de empleo cambia su pendiente rotando hacia la derecha, ya que aumenta el coste de la vacante al reducirse el tiempo que se espera que esté cubierta. 
\end{lnum}

En el nuevo equilibrio el desempleo aumenta, pero el efecto sobre las vacantes es indeterminado, al desplazarse las dos curvas. 

El equilibrio también dependerá de la tecnología de emparejamiento. Por lo que se podrían valorar los efectos de la creación de una Institución de Ayuda a la búsqueda de Empleo pública (por ejemplo, INE o SEPE)\footnote{%
    Por ejemplo, si se reduce la eficiencia en el emparejamiento el efecto será similar al de un aumento de la tasa de destrucción de empleo. En definitiva, los factores que generen rigideces en el mercado laboral provocarán un desplazamiento de la curva hacia fuera. Por el contrario, los factores que aumenten la eficiencia en la búsqueda o corrijan las posibles rigideces del mercado de trabajo, provocarán un desplazamiento hacia dentro de la curva.
} Cambios en el funcionamiento de las instituciones del mercado laboral también pueden desplazar la Curva de Beveridge.



En síntesis, los modelos de búsqueda ofrecen una explicación a la existencia de desempleo friccional como consecuencia del proceso costoso de emparejamiento de trabajadores y vacantes.\footnote{%
    Pero el modelo DMP también ha sido objeto de críticas. Por ejemplo, SHIMER (2005) lo ha acuestionado dado que no ajustaba bien las variaciones cíclicas que se producen en el paro y en las vacantes en Estados Unidos como consecuencia de las perturbaciones que se producen en los precios.\\También se han ido realizando extensiones incorporando, el crecimiento tecnológico (PISSARIDES, 2000) o la aversión al riesgo (ACEMOGLU y SHIMER, 1999).
} Sin embargo, el modelo no resulta suficiente para explicar el paro de larga duración y no consiguen explicar plenamente las fluctuaciones cíclicas del desempleo y las vacantes. 

En este sentido, la destrucción de empleos parece ser mucho más variable que la creación de puestos de trabajo ya que las reducciones del empleo en las recesiones se deben principalmente en la pérdida de puestos de trabajo existentes y obedecen en menor medida a disminuciones en el ritmo de creación de nuevos puestos de trabajo. Una posible explicación es la existencia de costes de ajuste en la dinámica de la demanda de trabajo, que sea analiza a continuación.

\subsection{Alternativas de especificación de la función de emparejamiento}


\clearpage
\section{Modelos de costes de ajuste} 
La teoría neoclásica analiza el mercado de trabajo en un contexto estático, sin considerar los aspectos temporales de las decisiones de los agentes en el mercado. Sin embargo, la introducción de un entorno dinámico implica que los agentes también consideran las consecuencias que se derivan de sus decisiones en periodos futuros. 
\begin{la}
    \item Por un lado, la oferta de trabajo intertemporal incluye el fenómeno de sustitución intertemporal de consumo y ocio, es decir, la posibilidad de sustituir consumo de bienes por ocio en el tiempo, distinguiendo entre cambios pennanentes y transitorios en el perfil temporal del salario [\authorfont{Ver Tema 3.A.25}].
    \item Por otro lado, la teoría dinámica de la demanda de trabajo estudia cómo se ajusta la demanda de trabajo a lo largo del tiempo tras una perturbación que haya alterado la demanda de factores en el equilibrio.
\end{la}

El análisis dinámico de la demanda de trabajo supone que cuando la empresa maximizadora de beneficios toma decisiones de demanda de trabajo en un contexto intertemporal la cuestión residirá en el ajuste óptimo de su plantilla, lo que resulta costoso.\footnote{%
    Por tanto, en el análisis dinámico de la demanda de trabajo es necesario la existencia de costes de ajuste en la cantidad de trabajo demandada por las empresas y las implicaciones que tienen estos costes de ajuste para el equilibrio del mercado de trabajo.
}

Podemos definir los costes de ajuste como:

Todos aquellos gastos en los que incurren las empresas según ajustan el tamaño de su fuerza de trabajo. 

Estos costes de ajuste son un elemento esencial en la modelización de la demanda de trabajo, pues la empresa a la hora de decidir modificar su plantilla tiene que tener en cuenta los costes de ajuste a los que tendrá que hacer frente.\footnote{%
    Es posible distinguir entre:
\begin{lnum}
    \item Coste internos: Se producen cuando la empresa lleva a cabo una reorganización del trabajo que genera una pérdida temporal de eficiencia. Por ejemplo, la adaptación de los trabajadores a nuevos procesos productivos o de los nuevos trabajadores a la empresa. Estos costes son dificiles de medir, ya que no se contabilizan de forma explícita por la empresa, por lo que suelen tener un carácter implícito.
    \item Costes externos: se producen cuando la empresa incurre en un gasto asociado a la contratación y despido de trabajadores. En general, tendrán un carácter más explícito, dado por las indemnizaciones por despido o por los costes de contratación (anuncio, selección, formación, etc).
    \item Costes brutos: son debidos a variaciones relativas a la entrada y salida de trabajadores, es decir a los cambios en el número total de trabajadores. 
    \item Coste netos: se incurre en ellos aunque no haya cambios en el empleo total, si no que se debe a la contratación de trabajadores sólo para reemplazar bajas.
\end{lnum}
} Como es dificil realizar una medición de los costes con mucha desagregación, se analizarán los costes variables (dependientes del número de trabajadores objeto de rotación en la empresa), y costes fijos (no dependen del número de trabajadores).

Los resultados alcanzados difieren tanto en función de cómo modelicemos los costes de ajuste, como en función de si la empresa se mueve en un marco de perfecta certidumbre o incertidumbre. Sin embargo, el marco general de los costes de ajuste se puede exprezsar de la siguiente forma:

\eqblock{
\begin{aligned}
    &\max{\pi}=\int_0^\infty e^{-rt}[f(L)-wL-C(\dot L)]\mathrm{d}t \qquad\text{Donde,}
    \quad \begin{cases}
        r\quad &\equiv \text{\scriptsize Tasa de descuento beneficios futuros}\\
        C(\dot L_t)\quad &\equiv \text{\scriptsize Costes de ajuste de plantilla}\\
        \dot L_t\quad &\equiv \text{\scriptsize Variación de la plantilla, $\Delta L_t$}\\
    \end{cases}\\[2em]
    L^*\quad &\to \quad \text{CPO:}\quad f'(L^*)=w+rb\dot L-b\dot L\qquad \text{Y, si $\exists\Delta L$, entonces:}\\[1em]
    C(\dot L)=
    &\begin{cases}
        C(\dot L)_i, \quad \text{si} \quad |\dot L|>0;\\
        0, \quad \text{si} \quad |\dot L|=0;
    \end{cases}
\end{aligned}
}{}

\subsection{Modelo de referencia: Costes de ajuste convexos o cuadráticos}
Empecemos por suponer que la Empresa enfrenta unos costes de ajuste convexos y ásimetricos (con pendiente mayor si se enfrenta una amortización de puestos de trabajo)\footnote{%
    Inicialmente se asumía, por su enorme simplicidad analítica, que dichos costes eran simétricos, de forma que daba igual que la empresa estuviese aumentando plantilla que reduciéndola. Sin embargo la evidencia parece apuntar en que en multitud de economías desarrolladas, especialmente europeas, los costes de reducción de plantilla eran superiores a los de aumento de plantilla, ya que la compensación para trabajadores desempleados era muy elevada.\\ En contraste, con datos para Estados Unidos, HAMERMESH (1993) comprobó que los costes de contratación podían ser incluso mayores que los de despido, algo que corroboran ABOWD y KRAMARZ (2003) para Francia en el caso de trabajadores cualificados.
} y una decisión intertemporal de demanda de trabajadores en un contexto de información completa. 

La existencia de costes de ajuste tiene las siguientes implicaciones en las decisiones de la empresa: 
\begin{lnum}
    \item La existencia de costes de ajuste hace que las decisiones tomadas en el periodo áctual extiendan sus efectos en los siguientes periodos y las decisiones tomadas en periodos anteriores condicionen las decisiones actuales.
    \item La empresa no ajusta su plantilla cada periodo, sino que suaviza la senda de ajuste de manera que el nivel de empleo existente en la empresa se acerca gradualmente a su nivel óptimo.
    \item La empresa anticipa reducciones futuras permanentes del salario real y empieza a ajustar su plantilla hacia el nivel óptimo antes de que dicha reducción se produzca.
    \item Si se producen perturbaciones transitorias que aumenten la productividad del trabajo, la empresa aumentará su demanda de trabajo significativamente en el periodo en el que se ha producido la perturbación e inmediatamente después la reducirá para iniciar el ajuste gradual de la plantilla hasta su nivel de partida.
\end{lnum}

\subsubsection{Costes variables}
En el gráfico representamos la curva de costes variables. 

\imagenfit{CAjuste_CV.png}{Curvas de costes variables de ajuste}{Apuntes ICEX-CECO. Tema 3.A.26}

La curva de costes variables incorpora el importante supuesto de que los costes de ajuste aumentan a tasa creciente, independientemente de que la empresa esté expandiendo o contrayendo su plantilla (como hemos analizado anteriormente, puesto que los costes marginales aumentan con la rotación).

\paragraph{Impicaciones: Costes de ajuste vs. Neoclásica}
La dinámica de la demanda de trabajo funcionaría de la siguiente manera: si suponemos que el output aumenta por una perturbación positiva, y las empresas esperan que este aumento sea permanente, aumentarán su demanda de empleo. Pero como es costoso realizar la transición inmediata a un nuevo equilibrio, la empresa procederá lentamente. Ocurrirá a la inversa si la perturbación es negativa, aunque los despidos se harán incluso más lentamente, puesto que los costes de eliminar empleos son estrictamente mayores que los de contratación.

\subsubsection{Costes fijos}
Éstos son independientes del número de trabajadores objeto de rotación, por lo que la empresa tiene dos opciones ante una perturbación: 
\begin{la}
    \item Permanecer con su nivel de empleo actual; o,
    \item Ajustar su demanda de trabajo.
\end{la} 

\paragraph{Implicaciones: Costes de Ajuste vs. Neoclásica}
En este contexto, la empresa podrá ajustar su fuerza de trabajo de forma inmediata, ya que incurre en los mismos costes que si lo hiciera de manera gradual. Las decisiones de ajuste de plantilla dependerán de qué alternativa le reporta mayores beneficios. Además, si los costes de ajuste son cuantiosos, las variaciones en el empleo serán repentinas y grandes. 

Siguiendo la clasificación de Nickell (1986), el supuesto de costes de ajuste cuadráticos que hemos analizado, cuyo supuesto es que los costes de contratación no se distinguen de los costes de despido, se enmarca en los costes de ajuste estrictamente convexos, que es el modelo teórico básico utilizado por numerosos autores desde los setenta. No obstante, existe un análisis similar que realiza el propio autor acerca de otros contextos, como el análisis de monopsonio dinámico o de costes estrictamente lineales,\footnote{%
    A la hora de realizar trabajo empirico, la mayor parte de autores opta por analizar modelos con costes estrictamente convexos. pero es importante tener en cuenta que si existen tramos importantes con costes lineales. asumir co nvexidad puede llevar a conclusiones erróneas sobre el proceso de ajuste.
} que desarrollaremos a continuación.

\subsection{Modelo alternativo: Costes de ajustes lineales}
Es importante tener en cuenta que al existir algunos tramos importantes con costes lineales, asumir convexidad puede llevar a conclusiones erróneas sobre el proceso de ajuste del empleo,\footnote{%
    Por ejemplo, este hecho no permite explicar los ajustes asimétricos y el empleo discontinuo. NICKELL (1986) asume que pueden existir tramos no convexos y, por tanto, la función de costes puede ser lineal en un tramo considerable,
} por lo tanto conviene ver la formulación lineal, que se utiliza comúnmente en la actualidad. La función es muy similar a la del caso de costes estrictamente convexos:\footnote{%
    La representación que aparece a continuación es la que proporciona NICKELL (1986) para explicar este tipo de costes. nos mostrará cómo a veces las empresas contratan, despiden, o se quedan como están.
}

\imagenfit{CAjuste_Lineales.png}{Función de Ajustes lineales}{Apuntes ICEX-CECO. Tema 3.A.26}

La empresa ajustará su nivel de trabajadores lh (o, respectivamente, L1 ) si esto es más alto (más bajo) que el nivel inicial de empleo L0 . Si no, significa que L0 se mantiene entre Lh y Lr, su nivel óptimo:

$$
\begin{aligned}
    &\begin{aligned}
        \text{(1)}\qquad f'(L_h)=w+rc_h\\
        \text{(2)}\qquad f'(L_f)=w+rc_f\\
    \end{aligned}\\[2em]
    L=&
    \left\{\begin{aligned}
        L_h,\quad \text{si}\quad &L_0\leq L_h; \quad &\text{\scriptsize (contratar)}\\
        L_0,\quad \text{si}\quad &L_h\leq L_0\leq L_f;\quad &\text{\scriptsize (mantener)}\\
        L_f,\quad \text{si}\quad &L_f\leq L_0; \quad &\text{\scriptsize (despedir)}
    \end{aligned}\right.
\end{aligned}$$


Los costes de despido y contratación tienen efectos contrapuestos en la demanda de trabajo, como se observa en la siguiente figura: 
\begin{lnum}
    \item Un incremento en el coste de despido dificulta a la empresa continuar reduciendo personal pero no crea ningún incentivo para contratar.
    \item Una reducción en los costes de contratación tiene un efecto positivo en el empleo, dado que incrementa el valor $L_h$ de empleo óptimo.
\end{lnum}

\imagenfit{CostesLineales-Ajuste.png}{Demanda de factor productivo trabajo $L$ y Ajustes dinámicos del factor, en función del nivel inicial}{Apuntes ICEX-CECO. Tema 3.A.26}

\subsection{Implicaciones}
Las principales implicaciones que se pueden derivar de este modelo de dinámica de la demanda de trabajo es que el manejo adecuado de los costes de contratación y despido garantiza una mayor estabilidad en la demanda de mano de obra. Estos costes de ajuste se pueden cuantificar y comparar. Por ejemplo, en Estados Unidos los costes de contratación suelen ser más altos que los de despido, mientras que en Francia pasa lo contrario. Esto lleva a implicaciones importantes, que pone de manifiesto CAHUC et al. (2014):\footnote{%
    Para concluir, comparando ambos países, los autores sostienen que (1) las tasas de desempleo de Francia suelen ser mayores (por la poca flexibilidad al despido); (2) en igualdad de niveles salariales, la producción será mayor en Estados Unidos; (3) pese a que en situaciones expansivas es más difícil contratar empleados en este país, y (4) en situaciones recesivas es más difícil despedir en Francia.
}
\begin{lnum}
    \item Un aumento en los costes de despido impide el ajuste a la baja del personal, y elimina los incentivos para contratar. Un aumento de los costes de contratación impide nuevas contrataciones.
    \item Un coste menor de contratación siempre tiene efectos positivos en el empleo ya que aumenta el nivel óptimo.
\end{lnum}  

Además, el ajuste estudiado puede ser diferente dependiendo de la composición de la mano de obra. Por ejemplo, en el caso de empleados de baja cualificación, los tiempos de ajuste suelen ser más cortos. 

Algunos trabajos publicados en los años 80 y principios de los 90 han tratado de estimar también la velocidad del ajuste de la demanda de trabajo a través de la adopción de representaciones cuadráticas y simétricas de los costes de ajuste (LAZEAR, 1990, ABRAHAM y HOUSEMAN, 1993). Se ha observado que si las empresas se ajustan más rápidamente, los desequilibrios entre el empleo deseado y el empleo existente son más acusados. Además, se ha demostrado que Estados Unidos es el país occidental donde más rápido se ajusta el empleo. También se ha observado que el grado de sindicalización no tiene por qué ser un factor determinante de la velocidad de ajuste.

\clearpage
\section{Conclusión}

\subsection{Recapitulación}

\subsection{Extensiones y relación con otras partes del Temario}

\subsection{Opinión}

\subsubsection{Idea final (Salida o cierra)}



\clearpage
\pagenumbering{gobble}
\MyBibliography
\MyBibItem{Autor}{Título}{Año}{DOI -- LINK}


% --- Anexos (no aparecen en el índice) ---
\clearpage
\anexo{\textbf{Preguntas Test}}
%\input{\TemaCode/questions.tex} 


\end{document}